% Draft section: the scale-rich construction as a nested hierarchy.
% Numbers: rg/NOTES.md and rg/results/*.json (scripts in rg/).
\section{The splitting hierarchy and a generation flow}
\label{sec:rg}

\paragraph{Cells as nested complexes.}
In the scale-rich (SR) construction every ligament grows from a uniform point,
at a uniform angle, until it meets the boundary of the convex cell holding
that point. Each step therefore cuts one convex cell into two, and the material
is a binary hierarchy of cells. A cell is a complex whose members are its two
children and the segments of its chord. Its boundary vertices are the
junctions that interior chords make with its sides. Ligament $t$ has thickness
$\lambda_t=\lambda_0 t^{-\alpha}$, which labels the generation; thickness does
not change the topology.

\paragraph{The hierarchy is a random binary search tree.}
Let $x$ be uniform in a convex body $K$ and let $\theta$ be independent of
$x$. For fixed $\theta$ the offset $s$ of $x$ perpendicular to $\theta$ has
density $\ell(s)/|K|$, with $\ell(s)$ the chord length, and the area fraction
on one side of the chord is $F(s)=\int_{-\infty}^{s}\ell/|K|$, the distribution
function of $s$. Hence
\begin{equation}
  r \equiv \frac{|K_0|}{|K|} \sim \mathrm{Uniform}(0,1)
\end{equation}
whatever the shape of $K$ and the law of $\theta$. Cells are chosen with
probability proportional to area, so the area-labelled splitting tree has
exactly the law of a binary search tree built from uniform keys. The $T+1$
cell areas are uniform spacings, $\sum A^2 = 2/(T+2)$ and
$P(TA>x)\to e^{-x}$. The cell split at step $t$ sits at depth
$2\ln t$ with variance $2\ln t$ (measured slopes $1.983$ and $1.980$). Area
moments scale ordinarily, $\sum A^q\sim t^{1-q}$, so the multiscaling of
Krapivsky and Ben-Naim is confined to lengths in the rectangular model
(reproduced: $\gamma(a,0)=(a-\sqrt{a^2+4})/2$ to 2\%). It is absent for
isotropic angles ($\sum P^a\sim t^{1-a/2}$ within $0.013$).

\paragraph{The coupling between generations is exactly marginal.}
A chord ends on existing ligaments with probability proportional to length,
and $\langle\ell_h\rangle\simeq c_\ell h^{-1/2}$ with $c_\ell=1.161$. The
birth index of the host is therefore $h=tU^2$ with $U$ uniform, and the
thickness ratio between a chord and its host,
\begin{equation}
  \frac{\lambda_t}{\lambda_h} = U^{2\alpha},\qquad
  \Big\langle \Big(\frac{h}{t}\Big)^{\alpha}\Big\rangle=\frac{1}{1+2\alpha},
\end{equation}
does not depend on $t$ (measured $0.667, 0.498, 0.330$ at $\alpha=\tfrac14,\tfrac12,1$).
Under coarse-graining in $\ln t$, $\alpha$ is thus relevant for the density,
$\rho(t)\propto t^{1/2-\alpha}$, which is the jamming line at $\alpha=\tfrac12$.
It is exactly marginal for the coupling of each generation to the frame it
hangs on. Moments $\langle(\lambda_h/\lambda_t)^p\rangle=1/(1-2p\alpha)$
diverge at $\alpha=1/(2p)$. For host bending ($p=3$) this happens at
$\alpha=1/6$.

\paragraph{Generation flow of the conductance.}
Give each segment the conductance $\lambda/\ell_{\rm seg}$. Adding chord $t$,
of conductance $g_t=\lambda_t/\ell_t$, between points $P$ and $Q$ raises the
dissipated power at unit field by the rank-one amount
\begin{equation}
  \Delta P_t=\frac{g_t\,\Delta V_t^2}{1+g_tR_t},
  \qquad
  \sigma(T)=\sum_{t\le T}\psi_t\,\lambda_t\ell_t\cos^2\theta_t ,
\end{equation}
where $\Delta V_t$ is the pre-existing potential difference and $R_t$ the
$PQ$ resistance. This is exact (checked to $4\times10^{-7}$). For
$\alpha<\tfrac12$ the increments $\psi_t$ are stationary in $\ln t$. The ratio
$\Phi=\sigma/\sigma_{\rm aff}$ to the affine bound
$\sigma_{\rm aff}=\sum\lambda\ell\cos^2\theta=\rho/2$ therefore converges to a
line of fixed points labelled by $\alpha$:
$\Phi=0.59,\,0.70,\,0.82$ at $\alpha=0,\,0.2,\,0.4$ ($\Phi(T)$ at $T=3\times10^4$:
$0.597,\,0.691,\,0.757$; the remaining transient decays as $T^{-(1/2-\alpha)}$).
For $\alpha>\tfrac12$ the late increments exceed one. The thick early wires are
nearly equipotential and focus the field into the gaps that late chords
bridge. A hierarchical surrogate that smears each chord into its cell and
laminates the children across it (bond moving in continuum form) is exact on
its own hierarchy. It reproduces the stationarity and the trend in $\alpha$,
but lies about $0.2$ below the network ($\Phi_{\rm lam}=0.37,\,0.48,\,0.62$).
The density exponent of transport is $m=1$ for every $\alpha$.

\paragraph{Elasticity.}
With pinned joints the SR network is floppy: junctions are T-shaped, with
degree three and collinear host segments. Its stiffness is therefore
bending-stabilised. In an Euler--Bernoulli frame at $\rho=0.4$ the
per-ligament energy decays with birth index as $i^{-x}$ with
$x=0.78,\,0.83,\,1.04$ at $\alpha=0,\,0.3,\,0.6$, close to the reported
$\tfrac34+\tfrac\alpha2$. The split of this law into degree and thickness
exponents cannot be identified within one ensemble, because both are functions
of $i$. Scaling all thicknesses bounds the density exponent,
$1\le m\le 3$. Growing $T$ at fixed $\lambda_0$ gives $m_T\simeq1.7$--$2.5$ for
$0\le\alpha\le0.4$ with no divergence as $\alpha\to\tfrac12$. Two couplings
flow: the slenderness of the newest generation, $s_t\rho(t)\to2c_\ell^2/(1-2\alpha)$,
and the marginal host ratio $U^{-2\alpha}$. Increments do not collapse on
$s_t$ alone.

\paragraph{What the hierarchy buys.}
The exact contraction of the cell tree (Dirichlet-to-Neumann messages,
chords as separators) is planar nested dissection. Its root width grows as
$\sqrt T$ ($w\simeq5n^{0.53}$ for a cell holding $n$ chords), so it gives no
computational gain. Chords end on the sides of ancestors, so a cell is not a
two-terminal bond and admits no closed scalar recursion. What the hierarchy
does give is structural and exact: the binary-search-tree law of the cells,
the stationary host law, and through them a generation flow in which $\alpha$
labels a line of fixed points.
